\section{Reading the statements together without inventing history}
\label{sec:statement-context-relations}

The balance sheet gives a dated description of the company. The income,
cash-flow and equity statements explain different aspects of the movement
between dates. Reading them together requires more than adding amounts with
similar labels. An equity balance is measured at an instant, profit accrues
over a period, a cash-flow adjustment can concern a different population of
items from the balance-sheet caption, and a subtotal calculated by the
analyst supplies no independent observation against which to check itself.

The retained Apple statement networks already record those relationships
and unresolved differences. The complete context draft makes their roles
explicit in the language. Its combined view preserves 256 source observations
and all 143 original relation records, including their original explanations
and diagnostic classifications. It gives 135 displayed differences, six
derived subtotals and two calculations that remain unresolved because an
operand is a reported dash. This accounts for the full retained network;
it does not establish the missing underlying company history.

\subsection{A reporting period can contain several balance dates}

Apple's annual equity statement is a report for fiscal 2024, but its opening
equity is measured at 30 September 2023 and its closing equity at
28 September 2024 \citep[p.~32]{apple2024}. The portable observation retains
the annual report-period string. The original pointed reading separately
records the measurement date. The adapter follows that pointer and retains
both meanings: the date of the stock and the period of the statement in
which it appears. It does not overwrite one with the other.

The same distinction matters for cash. An opening cash amount is a stock;
operating cash flow is a movement during the period. A cash bridge joins
the two endpoints and that intervening movement. A snapshot subtotal instead
compares amounts measured at the same instant. The language therefore
declares the relation's purpose as well as its operands. Same-period income
comparisons, same-date equity comparisons, stock-flow bridges, net-asset
changes and explanatory differences have different context rules.

For example, the displayed annual change in net assets is compared with the
displayed change in equity, using both sets of reported endpoints. With
\(A_t\), \(L_t\) and \(E_t\) denoting assets, liabilities and equity in
millions of dollars at the closing date, and subscript \(0\) denoting
the opening date, the comparison is
\begin{align}
\label{eq:statement-context-equity-change}
 (A_t-A_0)-(L_t-L_0)
 &= (364{,}980-352{,}583)-(308{,}030-290{,}437)\notag\\
 &=12{,}397-17{,}593=-5{,}196\notag\\
 &=56{,}950-62{,}146=E_t-E_0.
\end{align}
All operands come from the retained balance-sheet readings
\citep[p.~31]{apple2024}. Recomputing the right side from the equity
endpoints gives it a distinct observed basis. The corresponding interim
comparison gives 9,808. These displayed reconciliations retain the original
precision and consolidated scope; they do not identify the transactions
that produced the changes.

An annual-to-interim endpoint comparison also retains both source contexts.
The annual closing equity and the next quarter's opening equity are
reported at the same date, even though they appear in different reports.
Their equality can be compared in the combined named-source view. An
interim-only view cannot use the unselected annual observation. Its five
cross-period comparisons therefore remain unresolved, with their original
relation records still visible. No later source is silently supplied to
complete them, and the annual and interim bases are not generally declared
interchangeable.

\subsection{A subtotal needs an independent target to be a check}

The source audit exposed an evidentiary weakness in the earlier annual
relationship ledger. Three analyst-calculated totals---other equity
movements, operating reconciliation adjustments and operating asset/liability
adjustments---were stored as fixed numeric targets and labelled displayed
pass/fail checks. Reproducing those totals checks the calculation against
its saved value. It does not supply an independent disclosed total or an
additional accounting identity. The interim ledger already treated the
corresponding calculations as explanatory subtotals.

The new interpretation classifies all six as derived quantities and returns
no independent residual for them. Their original constants, coefficients
and classifications remain in the retained records, so the reason for the
change can be examined. The source amounts and earlier arithmetic are not
rewritten. The correction concerns the strength of the evidence attributed
to a calculated subtotal.

Other differences have a different purpose. The annual opening cash-flow
population exceeds balance-sheet cash by 772. The annual payable comparison
leaves 329, and interim tax expense less cash tax paid leaves
\(-12{,}397\). Those differences remain explanations to investigate, not
amounts to post into a balancing account. Even the zero vendor-receivable
difference does not identify a collection schedule, gross transactions or
their dates. The full-context result retains each relation's populations
and declared diagnostic purpose beside its arithmetic.

Two annual calculations cannot be completed under the new baseline because
the debt-issuance operand is a dash. The earlier displayed-total calculation
used an explicit no-contribution convention. That convention and its result
remain in the historical evidence, while the new source interpretation
preserves the dash and leaves the financing total and term-debt diagnostic
unresolved. Using such a convention again would require its own explicit
premise. Similarly, five interim relations retain the annual concepts absent
from the interim table. Summing the rows that are actually listed does not
establish that an omitted class of transaction was zero.

\subsection{Local compatibility is not a company history}

The optional rounding intervals in Section~\ref{sec:company-source-state}
make it possible to ask a further, conditional question: could an individual
measurement identity hold for values inside those intervals? This requires
two separate premises. One specifies the rounding rule; the other asserts
that the proposed identity concerns the same underlying measurements.
A zero displayed residual alone supplies neither premise.

Let a proposed relation have residual \(r=\sum_i a_i x_i\), including
the target with coefficient \(-1\). Before computing an interval for
\(r\), collect repeated occurrences of the same observation and add their
coefficients. Thus \(x-x=0\) exactly, rather than the difference of two
independently varying copies of \(x\). For a positive coefficient, the
smallest contribution uses the lower endpoint and the largest uses the
upper. For a negative coefficient those choices reverse. Summing the
contributions gives the residual endpoints; an endpoint belongs to the
result only if every nonzero contribution can attain its relevant endpoint.

For the annual current-assets relation, six positive displayed components
sum to the displayed total of 152,987. Under the proposed nearest-rounding
rule with quantum one million, write the component errors as
\(e_1,\ldots,e_6\) and the total's error as \(e_T\). Every one of these
positive displayed measurements has error in \([-1/2,1/2)\). The residual
of the underlying measurements is therefore
\begin{equation}
\label{eq:statement-context-residual-range}
  r=\sum_{i=1}^{6}e_i-e_T\in(-7/2,\,7/2).
\end{equation}
The lower extreme would require the total's excluded upper endpoint, and
the upper extreme would require the components' excluded upper endpoints.
Zero lies inside this range. Under the separately proposed identity, that
establishes local compatibility with the intervals. It leaves the rest of
the statement network untested.

A simple counterexample shows why this distinction is necessary. If
\(x\in[0,1]\), the proposed equality \(x=0\) is locally compatible,
and so is \(x=1\). No single value satisfies both. A collection of local
interval tests is therefore not a witness for a joint feasible history.
The displayed-statement reference retains joint feasibility as uncomputed
and actual history as unknown. Section~\ref{sec:joint-source-measurements}
adds a separate conditional assessment of selected balance measurements.
That assessment creates neither legal books nor the claims, lots, contracts
and behavioral assumptions required to continue the company.

The complete observation and relation declarations have matched typed and
named-text forms in \path{RetainedStatementContextDraft/1}. Each operand
has its source identity, date or interval, unit and explicit alias. Relation
roles, signs, targets, source omissions and retained dependencies are part of
the declaration. Source binding rejects changed dates, units, populations,
aliases and roles; result verification recomputes the complete description.
The reference preserves per-share values and weighted-average share counts
alongside monetary amounts without summing unlike units.

The checked records are under \path{data/accounting-master/}:
\begin{quote}\small
\path{source-relations-20261024}\\
\path{statement-context-coverage-20261024}
\end{quote}
The former contains conditional interval diagnostics for the balance
relations; the latter supplies complete displayed-statement contexts and
explicit unresolved results. Together they make the source layer more
precise. General source-to-state mappings, shared method requirements,
admissible latent histories and independent acceptance remain subsequent
design obligations.
